\par\addvspace{7pt}\Needspace{4\baselineskip}\noindent\textbf{作业题}\par
\begin{exercise}\label{mit:6042-ps4-q2}
\textbf{6.042J，Spring 2015，PS4 第 2 题。}
设 $L$ 为标签集。带标签满二叉树递归定义为：单个带标签叶 $\langle\ell,\mathrm{leaf}\rangle$ 是树；若 $B,C$ 是树，则 $\langle\ell,B,C\rangle$ 是树。叶标签集和内部标签集也按两棵子树的对应集合取并，其中新根标签加入内部标签集。
“标签唯一”指两棵子树各自标签唯一、彼此的全部标签集不交，且新根标签不在任何子树中。设 $n_B,f_B$ 分别为不同内部标签和不同叶标签的数目。用结构归纳证明标签唯一时 $f_B=n_B+1$，并指出哪里使用了唯一性。
\end{exercise}
\sourceof{PS4，物理第 1--2 页；编者独立解答。}
\begin{solution}
基例只有一个叶标签，$f=1,n=0$。构造步设 $B=\langle\ell,C,D\rangle$，归纳假设为 $f_C=n_C+1$、$f_D=n_D+1$。由于两个子树的标签集不交，$f_B=f_C+f_D$；又由于新根标签不在子树中且子树内部标签不交，$n_B=n_C+n_D+1$。于是 $f_B=n_C+n_D+2=n_B+1$。这里两条加法等式正是使用标签唯一性的地方。若允许两个叶同用一个标签，根用另一个标签，则 $f=1,n=1$，结论失效。
\end{solution}

\begin{exercise}\label{mit:6042-ps6-q2}
\textbf{6.042J，Spring 2015，PS6 第 2 题。}
\begin{enumerate}[label=(\alph*)]
\item 构造有向图和不同顶点 $u,v$，使 $u$ 可达 $v$ 且 $v$ 可达 $u$，但没有有向圈同时包含二者。
\item 证明：若存在从 $v$ 出发又回到 $v$ 的正长度有向游走，则存在包含 $v$ 的有向圈。
\end{enumerate}
\end{exercise}
\sourceof{PS6，物理第 1 页；编者独立解答。}
\begin{solution}
(a) 取顶点 $u,w,v$ 和弧 $u\to w,w\to u,v\to w,w\to v$。两方向都有长 $2$ 路径。任何同时包含 $u,v$ 的闭合游走都必须两次经过 $w$，因此不是简单有向圈。

(b) 在所有经过 $v$ 的正长度闭合游走中选最短者 $W$。若内部顶点 $x$ 重复，删除两次出现 $x$ 之间的闭合段，仍得到经过 $v$ 的更短正长度闭合游走；若内部再次出现 $v$，取此前的闭合前缀也更短。因此除首尾 $v$ 外没有重复顶点，$W$ 就是有向圈。若允许自环，长 $1$ 的情形直接成立。
\end{solution}

\begin{exercise}\label{mit:6042-ps6-q3}
\textbf{6.042J，Spring 2015，PS6 第 3 题。}
竞赛图中每对不同顶点之间恰有一条定向弧。若 $u$ 到每个其他顶点的距离至多 $2$，称 $u$ 为王。
\begin{enumerate}[label=(\alph*)]
\item 构造 $10$ 顶点竞赛图，其中有一个出度为 $1$ 的王。
\item 构造 $5$ 顶点竞赛图，其中所有顶点都是王。
\item 证明每个出度最大的顶点都是王。
\end{enumerate}
\end{exercise}
\sourceof{PS6，物理第 1--2 页；编者独立解答。}
\begin{solution}
(a) 顶点为 $u,w,x_1,\ldots,x_8$。规定 $u\to w$，$w\to x_i$，$x_i\to u$，并按 $i<j$ 定向 $x_i\to x_j$。$u$ 唯一出邻点为 $w$，其余顶点均通过 $w$ 在两步内到达。

(b) 用模 $5$ 的顶点 $0,1,2,3,4$，规定 $i\to i+1$ 和 $i\to i+2$。差为 $1,2$ 的顶点直接可达，差为 $3$ 可用 $i\to i+1\to i+3$，差为 $4$ 可用 $i\to i+2\to i+4$。

(c) 设 $u$ 出度最大，$A=N^+(u)$。若某 $v$ 不能在两步内从 $u$ 到达，则 $v\to u$ 且对每个 $a\in A$ 都有 $v\to a$，否则 $u\to a\to v$。所以 $d^+(v)\ge |A|+1>d^+(u)$，矛盾。$n=1$ 时结论为空真。
\end{solution}

\begin{exercise}\label{mit:6042-ps7-q3}
\textbf{6.042J，Spring 2015，PS7 第 3 题。}
原图的四个简单图均用顶点 $1,\ldots,10$ 标号。为完全确定原图，下面列出边集；$12$ 表示无向边 $\{1,2\}$，$1(10)$ 表示 $\{1,10\}$。四图共同含五圈 $12,23,34,45,51$（$G_4$ 除外）：
\[
\begin{array}{c|l}
G_1&12,23,34,45,51,16,29,37,4(10),58,67,78,89,9(10),(10)6\\
G_2&12,23,34,45,51,16,27,38,4(10),59,67,78,8(10),(10)9,96\\
G_3&12,23,34,45,51,16,27,38,4(10),59,6(10),68,97,9(10),78,(10)8\\
G_4&12,23,34,45,56,67,78,89,91,1(10),4(10),7(10),95,26,83
\end{array}
\]
判定哪些图同构；对每对同构图给出一个同构，对每对不同构图给出区分它们的同构不变量，并至少证明其中一个性质确实为同构不变量。
\end{exercise}
\sourceof{PS7，物理第 1--2 页，图 1；编者独立解答，边集由原图核对。}
\begin{solution}
只有 $G_1,G_4$ 同构。按 $G_1$ 顶点 $1,\ldots,10$ 的顺序，其像可取
\[
(1,2,3,4,10,9,8,7,6,5).
\]
这是双射，逐一把上表 $G_1$ 的 $15$ 条边代入便得到 $G_4$ 的全部 $15$ 条边，故是同构。

$G_1,G_2,G_4$ 各有 $15$ 条边且所有顶点度数为 $3$；$G_3$ 有 $16$ 条边，顶点 $8,10$ 的度数为 $4$。边数已区分 $G_3$ 与其余三图。$G_2$ 有四圈 $1,2,7,6,1$，$G_1$ 没有四圈，故 $G_2$ 与 $G_1,G_4$ 均不同构。核对 $G_1$ 无四圈的方法是：其外侧和内侧分别是五圈，内外之间是一一对应的五条边。四圈必须使用偶数条内外边，且不能用四条（内外边构成匹配）；若用两条，则对应的一对外侧顶点和一对内侧顶点都须相邻。上表显示五条外侧邻接对应的内侧顶点分别为 $6,9$；$9,7$；$7,10$；$10,8$；$8,6$，均不相邻。因此无四圈。

证明圈长是同构不变量：若 $x_1,\ldots,x_k,x_1$ 是圈，同构 $f$ 保持相邻且保持各顶点互异，故 $f(x_1),\ldots,f(x_k),f(x_1)$ 也是同长圈；对 $f^{-1}$ 同理。因此“存在四圈”在同构两侧等价。
\end{solution}

\begin{exercise}\label{mit:6042-ps7-q4}
\textbf{6.042J，Spring 2015，PS7 第 4 题。}
称恰有两个度 $1$ 顶点、其他顶点度 $2$ 的简单图为“两端图”。原题把可沿一个序列排列、且仅相邻位置连边的图称为 line graph；此处按其实际定义译作\textbf{路径图}，以免与通常的线图概念混淆。
\begin{enumerate}[label=(\alph*)]
\item 反驳“每个两端图都是路径图”。
\item 找出以下归纳伪证的首个缺口：单边情形成立；假设每个 $n$ 边两端图都是路径图。取一个这样的 $G_n$，新增一条边并保持两端性质只能把新顶点连到端点，故所得 $G_{n+1}$ 是路径图，于是所有 $n+1$ 边两端图都成立。
\end{enumerate}
\end{exercise}
\sourceof{PS7，物理第 1--3 页；编者独立解答，术语按题内定义修订。}
\begin{solution}
(a) $K_2\sqcup C_3$ 是反例：单边分支给两个度 $1$ 顶点，三圈上的三个顶点度 $2$，却不连通，故不是路径图。

(b) 增添边的局部判断没有覆盖全部 $n+1$ 边两端图。把一个 $n$ 边两端图扩成某些 $n+1$ 边两端图，不能推出\textbf{任意} $n+1$ 边两端图都能如此产生。例如上述 $4$ 边反例中，删掉单边会留下度 $0$ 顶点，删掉圈边则产生另外两个度 $1$ 顶点；无论怎样删边都不能得到 $3$ 边两端图。首个无根据的推广就在把“我们构造出来的图”当作“所有更大图”。若补上连通条件，则沿端点走、每个内部点只有一个尚未使用的下一条边，会走遍全部顶点，正确结论确为路径图。
\end{solution}

\begin{exercise}\label{mit:6042-ps8-q1}
\textbf{6.042J，Spring 2015，PS8 第 1 题。}
证明有限\textbf{连通}带权图的边权两两不同时，其最小生成树唯一。原题未明说连通；这里补上保证生成树存在的条件。
\end{exercise}
\sourceof{PS8，物理第 1 页；编者独立解答。}
\begin{solution}
若 $M,N$ 是不同的最小生成树，选对称差 $M\triangle N$ 中权最小的边 $e$，交换名称使 $e\in M\setminus N$。在 $N+e$ 的唯一圈上，至少有一边 $f\notin M$，否则 $M$ 已含这个圈。因此 $f\in N\setminus M$，且不同权给 $w(f)>w(e)$。删去 $f$ 得到生成树 $N+e-f$，其权比 $N$ 小，矛盾。
\end{solution}

\begin{exercise}\label{mit:6042-ps8-q2}
\textbf{6.042J，Spring 2015，PS8 第 2 题。}
原图是如下 $11$ 顶点图：下标按模 $5$ 计算，顶点为 $v_0,\ldots,v_4,u_0,\ldots,u_4,w$；边为
\[
v_iv_{i+1},\qquad u_iv_{i-1},\qquad u_iv_{i+1},\qquad wu_i\quad(0\le i<5).
\]
它没有三角形。（a）给出 $4$ 染色；（b）证明不能 $3$ 染色。说明为什么高色数并不必然意味着包含大完全子图。
\end{exercise}
\sourceof{PS8，物理第 1 页，图 1；编者独立解答，原图按五圈 Mycielski 构造重标号。}
\begin{solution}
(a) 给 $v_0,\ldots,v_4$ 的颜色依次为 $1,2,1,2,3$；令 $u_i$ 与 $v_i$ 同色，并给 $w$ 色 $4$。五圈正确染色；$u_i$ 的两个原顶点邻居与 $v_i$ 相邻，所以与 $u_i$ 异色；$w$ 的邻居也都异于色 $4$。

(b) 若有三色染色，重命名使 $w$ 为色 $3$，则全部 $u_i$ 用色 $1,2$。把每个原顶点 $v_i$ 若为色 $3$ 就改成 $u_i$ 的颜色，其余不改。这给五圈一个正确二染色：一条 $v_iv_j$ 边不可能原来两端都是 $3$；若只有 $v_i$ 原为 $3$，则 $u_i$ 与 $v_j$ 相邻，换后仍异色。五圈为奇圈，不可能二染色，矛盾。

三角形若含 $w$，其余两点须是互不相邻的 $u$ 顶点，不可能；若不含 $w$，最多有一个 $u_i$，而它的两个原顶点邻居在五圈上不相邻。因此图无三角形，当然也不含 $K_r$（$r\ge3$），却有色数 $4$。
\end{solution}

\begin{exercise}\label{mit:6042-ps8-q3}
\textbf{6.042J，Spring 2015，PS8 第 3 题。}
四男四女的偏好如下，横线是尚未指定的两个人，任意补成严格排序：
\[
\begin{array}{c|cccc@{\qquad}c|cccc}
B_1&G_1&G_2&-&-&G_1&B_2&B_1&-&-\\
B_2&G_2&G_1&-&-&G_2&B_1&B_2&-&-\\
B_3&-&-&G_4&G_3&G_3&-&-&B_3&B_4\\
B_4&-&-&G_3&G_4&G_4&-&-&B_4&B_3
\end{array}
\]
\begin{enumerate}[label=(\alph*)]
\item 证明 $(B_1,G_1),(B_2,G_2),(B_3,G_3),(B_4,G_4)$ 在所有补全下都稳定。
\item 说明它既不是男方最优也不是男方最差匹配，因而不是任一侧申请的延迟接受算法输出。
\item 构造 $n$ 男 $n$ 女的偏好，使至少有 $2^{n/2}$ 个稳定匹配。原提示将双方各分成 $n/2$ 对，故隐含 $n$ 为偶数；先证明该情形，再讨论奇数。
\end{enumerate}
\end{exercise}
\sourceof{PS8，物理第 1--2 页；编者独立解答，(c) 的奇数条件另作说明。}
\begin{solution}
(a) 在第一组 $\{B_1,B_2;G_1,G_2\}$，两位男方都得到首选；第二组 $\{B_3,B_4;G_3,G_4\}$，两位女方都得到本组首选。因此各组内部没有阻塞对。跨组时，若男方来自第一组，他更喜欢本组伴侣；若女方来自第一组，她更喜欢本组伴侣。故任意补全下都没有跨组阻塞。

(b) 只交换第一组的伴侣，第一组女方均得到首选，仍稳定；此时 $B_1,B_2$ 比原匹配更差，所以原匹配不是男方最差。只交换第二组的伴侣，第二组男方都从第四选择提升为第三选择，仍稳定，所以原匹配不是男方最优。跨组判断与 (a) 相同。引理~\ref{lem:6042-stable} 说明两侧申请输出恰是这两个极端。

(c) 当 $n=2k$，按同一顺序划分成 $k$ 个双人组。每个人都把所有较早组排在本组之前，所有较晚组排在本组之后。在本组，男方喜欢同编号女方胜过另一位，女方喜欢另一位男方胜过同编号。每组的直配和交叉配都稳定：前者使男方在本组最优，后者使女方在本组最优。各组任取一种，共 $2^k$ 种。任一跨组潜在配对，属于较早组的一方更喜欢本组伴侣，故不会阻塞。

原断言若包含 $n=1$ 则错误，因为此时只有 $1<\sqrt2$ 个匹配。对于奇数 $n\ge3$，仍能达到所需下界，但需要补充三人组：用模 $3$ 编号，男方 $B_i$ 对本组的排序为 $G_i,G_{i+1},G_{i+2}$，女方 $G_j$ 为 $B_{j+1},B_{j+2},B_j$。三个循环匹配 $B_i\mapsto G_{i+s}$（$s=0,1,2$）均稳定；$s=0$ 男方得首选，$s=2$ 女方得首选，$s=1$ 时男方唯一更喜欢的是 $G_i$，而 $G_i$ 更喜欢她的现任 $B_{i+2}$。将一个三人组和 $(n-3)/2$ 个双人组按前述组间排序组合，得到 $3\cdot2^{(n-3)/2}>2^{n/2}$ 个稳定匹配。这是编者对原偶数提示的补充。
\end{solution}
